\chapter{Conclusion and outlook}
\label{ch:conclusion}

\section{Summary}
I transferred the method of \citet{bian2025} to a window-ventilated classroom of our university, using only
open-source software for the simulations. The answers to the questions of \cref{ch:intro} are:
\begin{enumerate}
  \item The method works with OpenFOAM. The pipeline from the room geometry to the trained model runs from
        scripts.
  \item My \cotwo{} solver reproduces OpenFOAM's own \cotwo{} within about 2\,\% when the flow is saved every
        second. Saving it only every 10~s, as in the first dataset, caused differences of up to 45\,\% in fast
        settings; the corrected dataset fixes this for the first 150~s. The forecast error is 4.5 times
        (transition runs) to 10 times (direct runs) smaller than persistence, and the model trained on the
        corrected data is more accurate on the corrected test data.
  \item Running the cases on 4 cores is 2.7 times faster, and 450~s instead of 600~s would save another 25\,\%.
        A lattice Boltzmann solver on the GPU was neither faster nor close enough to OpenFOAM. Among the available
        software only ANSYS Fluent supports the same model fully on the GPU.
  \item Training only from the equations did not reach a useful accuracy for this room (velocity error about
        50\,\%).
  \item For the real sensor data, the model has to work with few measurement points, window states and all
        windows and doors (\cref{ch:sensors}).
\end{enumerate}

\section{Limitations}
\begin{itemize}
  \item After 150~s the \cotwo{} is computed on the mean flow. In strongly fluctuating settings this keeps too much
        \cotwo{} in the room; in the paper's metric the difference stays at a few per cent.
  \item The seating area is a placeholder until the real seating plan is known.
  \item The dataset covers only W1, W3 and W6 with fixed doors.
  \item The model has not yet been tested on real measurements.
\end{itemize}

\section{Outlook}
\todo{sensor data, all windows and doors, control of the windows as in the paper}
